\section{Discussion and Conclusion}

Over one-third of regulatory substance entries cannot be linked to a molecular structure; CAS numbers introduce ambiguity rather than resolving it; and neither embeddings nor LLMs can substitute for structural information. Full manual reconciliation of non-structure entries (an estimated 3,150~person-years for REACH alone\footnote{Worst-case pairwise reconciliation (comparing every pair; real reconciliation would cluster) of $n \approx 11{,}250$ non-structure entries needs $n(n-1)/2 \approx 6.3\times10^{7}$ comparisons; at \textasciitilde{}100/day and 200 working days/year, that is $\approx 3{,}150$ person-years.}) is infeasible.

These findings show that downstream harmonisation cannot substitute for change at the source: the different regulatory instruments that define and collect substance data (REACH, CLP, and the sector-specific lists built on them) must encode structure-based and group-level semantics consistently, not only the systems that integrate them afterwards.
Three recommendations follow: (1)~require structure-based identifiers (InChIKey; MInChI~\cite{Heller2015} for mixtures, or representative-constituent structures~\cite{Kutsarova2019} for UVCBs where a single structure does not apply) across regulatory instruments; (2)~mandate explicit representation of substance groups using formal ontologies such as ChemOnt, so that group membership is computed automatically from structure via ClassyFire rather than maintained as a static, manually curated list --- shown feasible here by the 304 SKOS crosswalk triples (C5); (3)~publish regulatory data as linked data~\cite{BernersLee2006LinkedData} with persistent URIs and SKOS taxonomy links, enabling FAIR-compliant~\cite{Wilkinson2016} integration under the Interoperable Europe Act~\cite{InteroperableEuropeAct2024}.
Future work should assess generality beyond ECHA/EU data (e.g.\ US EPA, UK HSE) and develop curated resolution workflows for the 37.5\% of entries that resist automated linking, guided by the taxonomy introduced here.

\textit{Limitations:} Scoped to 18 ECHA-based datasets; ClassyFire ground truth is itself algorithmically assigned, not expert-validated, and a manual sample shows most apparent LLM errors select a chemically valid alternative label rather than a hallucinated one; LLM evaluation uses a single commercial model (Claude Sonnet~4.6); 68\% of non-structure entries resist automated alignment.
\textit{Code and data:} \url{https://doi.org/10.5281/zenodo.21512303}~\cite{VanHaute2026Zenodo}.
